\newpage
% \color{white}
% \pagecolor{black}

\ifdefined\useChapters
\chapter{Contra fogo}
\else
\chapter{}
\fi

O ar tremia, denso e quente, como se o dia tivesse se inflamado de repente.  
Não era o sol — era ela.

Entre o zumbido distante de carros e o estalo do asfalto quente, Larissa vinha pela rua com o rosto duro, o cabelo azul faiscando em tons de labareda viva.  
Aquela visão misturava perigo e beleza de um jeito que confundia — algo entre medo e fascínio.

— Sério mesmo... em plena luz do dia? — murmurei, sem saber se falava pra mim ou pro universo.

Ela não respondeu.  
Parou a poucos metros, o olhar tão fixo que parecia atravessar o ar.

— Eu não sei como você fugiu de lá — disse, aproximando-se com passos firmes —, mas eu vim te buscar de volta.

O coração disparou.  
Eu não tinha chance se ela quisesse me machucar, e ela parecia bem-disposta a isso.  
A menos de um metro, Larissa me agarrou pela gola.

— Você não vai correr? Pensei que fosse covarde — disse, entre ironia e raiva.

— Correr não ia adiantar — respondi, tentando parecer calmo. — Acho que você não quer me matar. Se quisesse, já teria feito.

Ela estreitou os olhos, e a mão no meu braço começou a esquentar.  
O cheiro de tecido queimado e pele ardendo subiu rápido — como frango frito, grotesco.  
O reflexo veio antes do pensamento: empurrei-a, desesperado.

Larissa caiu de costas, apoiou-se nas mãos e riu — um riso breve, ofendido.

— Então é assim? — ela se levantou, e junto dela o ar se contraiu. — Achei que você fosse indefeso.  
Mas já que é assim que quer...

Uma bola de fogo se formou na palma dela, crescendo até ofuscar o rosto.  
Disparou.

Não tive tempo pra pensar; o corpo reagiu sozinho.  
Senti o ar se distorcer entre nós — e, num impulso, a energia telecinética que eu mal entendia desviou o ataque.  
A bola explodiu contra um muro, espalhando faíscas e fumaça.

A sensação era estranha: tocar sem encostar, mover o mundo pelo susto.

Larissa franziu o cenho.  
— O que é isso?  
— Você não era o cara das ilusões? — perguntou, dando um passo pra trás. — Da outra vez achei que era truque, mas assim, tão perto, vejo que não.

O fogo voltou a dançar entre os dedos dela.  
— Não é à toa que o velho te quer vivo.

— Por quê? — perguntei. — Você acredita nas ideias dele?  
— Isso não te interessa.  
— Ele acha que uns devem dominar outros. E você?  
Ela hesitou — mínima fração de segundo, mas suficiente para a chama oscilar.

— A única coisa que importa é que eu vou te levar amarrado — disse, forçando firmeza. — E dessa vez você não vai escapar.

Começou a lançar pequenas bolas de fogo, rápidas, cadenciadas.  
Algumas acertavam o chão, outras roçavam minha roupa, deixando manchas escuras e cheiro de couro queimado.  
A cada golpe, ela parecia tentar convencer a si mesma de que tinha razão.

— Ai, para com isso, Larissa! — gritei, desviando o máximo que podia. — Você não precisa obedecer aquele velho.  
— Obedecer? — ela cuspiu a palavra. — Ele não manda em mim!

O tom saiu tão alto que me assustou.  
Ela respirava pesado, o rosto rubro, as mãos tremendo.  
O fogo não era mais controle — era fúria.

— Então por que faz o que ele quer? — insisti. — É medo ou costume?

Ela lançou uma labareda que passou a centímetros do meu rosto.  
— Você fala demais, garoto.

O calor era insuportável; a pele ardia, os olhos lacrimejavam.  
Tentei reagir, concentrar, mas o medo me quebrava.  
Cada tentativa de controle vinha acompanhada de um estalo, um chiado interno, como se o mundo dissesse “ainda não”.

Larissa girou o pulso, criando um círculo de fogo ao meu redor.  
A chama se ergueu como uma parede viva.  
O ar rareou.  
O som se tornou grave, abafado.

— Não adianta — disse ela. — Já era.

Senti o corpo ceder.  
As mãos suavam, o coração batia no pescoço.  
Mas dentro do medo havia algo: uma lembrança do dia em que consegui mover o ar pela primeira vez, o instante em que a confiança veio antes do pensamento.  
Tentei de novo.

Fechei os olhos.  
O calor me envolveu, depois veio o vento, a pressão no peito — e, por um segundo, o silêncio.

Quando abri os olhos, o fogo sumira.  
O chão também.

— O que é isso?! — gritei.

— Não se debate, cara, senão cai — disse uma voz acima de mim.

Um braço me segurava firme pelos ombros; a paisagem passava rápida lá embaixo.  
O vento batia no rosto, frio e libertador.

Olhei pro lado.  
Era ele — o volitador.

— Eu não quero fugir da luta — protestei.  
— Luta? — ele riu. — Você estava virando churrasco.  
Dei uma olhada pra baixo e entendi.  
A rua era só fumaça e fuligem.

— Se eu não tivesse chegado por trás dela e dado uma pancada na nuca, você já era — disse, com o tom calmo de quem conta o placar de um jogo.
— Eu detesto admitir, mas... — Respirei fundo. — Você tem razão.

Ele deu um meio sorriso.  
— Normal. A maioria das pessoas que eu salvo prefere discutir primeiro.

O vento cortou a conversa.  
Lá embaixo, as labaredas minguavam.  
E, pela primeira vez, percebi que o medo não vinha do fogo — vinha de mim.
